\documentclass[11pt,a4paper]{article}
\usepackage[T1]{fontenc}
\usepackage[utf8]{inputenc}
\usepackage{amsmath,amssymb,xcolor}
\usepackage[margin=25mm]{geometry}
\usepackage[colorlinks=true,urlcolor=blue]{hyperref}
\definecolor{nexorablue}{HTML}{2459B5}
\setlength{\parindent}{0pt}
\setlength{\parskip}{7pt}
\setlength{\emergencystretch}{2em}
\renewcommand{\familydefault}{\sfdefault}
\begin{document}

{\small\bfseries\color{nexorablue} NEXORA CPGE}\par
{\small MPSI / NEX-ENS-01 / Corrigé}\par
{\LARGE\bfseries Deux filtres sur un ensemble de mesures}\par
\vspace{0.25em}\hrule\vspace{0.6em}

\section*{Ensembles de départ}
On a $A=\{0,3,6,9,12\}$ et $B=\{8,9,10,11,12\}$. L’intersection vaut $\{9,12\}$.

\section*{Différence et complémentaire}
Les éléments appartenant à un seul filtre sont $\{0,3,6,8,10,11\}$. Ceux qui n’appartiennent à aucun filtre sont $\{1,2,4,5,7\}$.


\end{document}
